\ifdefined\MyWebBook\else
\documentclass[../textbook.tex]{subfiles}
\begin{document}
\fi
\bookchapter{训练数据工程基础}{ch:data-engineering}

模型从经过采集、解析、筛选、变换和采样的数据中学习。两个系统即使采用同一架构和相同的原始文件，也可能因为样本边界、去重规则、标签掩码和混合比例不同而优化不同的目标。因此，\term{数据工程}{Data Engineering}{}既要保证处理链能够持续运行，也要解释每个训练词元为何进入数据集、代表什么分布，以及能否追溯到原始来源。

训练数据的生产从原始资料出发，经治理、去重、混合和分片形成可复现、可持续消费的数据版本。这条处理链适用于预训练和后训练。样本与词元的表示见\bookxref{ch:tokenization}；各类监督数据的具体要求分别见预训练、监督微调、偏好对齐和可验证推理相关章节。

\bookxref{ch:autoregressive}定义了条件概率与有效预测位置；数据系统通过样本切分、标签构造和采样决定这些位置的内容与分布。为追溯一次训练使用的数据，需要记录统计单位、样本结构、准入规则、处理版本与存储位置。

本章的主要符号见表\ref{tab:data-engineering-symbols}。概率按记录还是按词元计算，将在使用处单独声明；容量统一用词元位置计量。
\begin{table}[htbp]
\centering
\caption{主要符号与计量对象}
\label{tab:data-engineering-symbols}
\begin{tabular}{ll}
\toprule
符号 & 含义与范围\\
\midrule
$P_\star,P_0,Q$ & 目标、原始采集及最终训练分布\\
$A,B$ & 用于比较的有限片段集合\\
$J(A,B),\tau$ & Jaccard 相似度与确认阈值，取值 $[0,1]$\\
$m,b,r$ & 签名行数、带数、每带行数，$m=br$\\
$n_i,p_i,q_i$ & 来源规模、实际采样概率、目标风险权重\\
$\alpha$ & 来源混合指数，本章取 $[0,1]$\\
$L_i,M,T$ & 样本长度、总长度、单块容量，单位为词元位置\\
$K,R,P$ & 块数、丢弃尾长、填充位置数\\
\bottomrule
\end{tabular}
\end{table}

\section{训练数据的组织}

\subsection{样本依赖性}
设目标应用中的样本分布为 $P_\star$，原始采集分布为 $P_0$，处理链最终提供给训练器的分布为 $Q$。训练风险 $\mathbb E_{z\sim Q}[\ell_\theta(z)]$ 基于\emph{最终供给分布} $Q$，与目标应用中期望降低的\emph{真实分布} $P_\star$ 风险 $\mathbb E_{z\sim P_\star}[\ell_\theta(z)]$ 一般不同。网站可访问性、抓取频率、语言识别器和长度筛选都会改变 $Q$。数据量增加只能降低对当前采样过程的部分估计误差，分布偏移则依然存在。

\term{样本单位}{Sampling Unit}{}需要随任务明确：文档、对话、问题家族、用户或时间窗口都是不同单位。同一篇文档的多个切片、同一问题的不同改写、同一对话的多个轮次具有相关性。把它们当作独立样本，会夸大有效信息量，也会使随机划分产生泄漏。

考虑 $G$ 个相互独立的组，每组含 $m$ 个观测，组内观测方差均为 $\sigma^2$，任意两个不同观测的相关系数均为 $\rho$，总数 $N=Gm$。样本均值的方差为
\begin{align}
\operatorname{Var}(\bar Z)
 &=\frac{G[m\sigma^2+m(m-1)\rho\sigma^2]}{G^2m^2}\\
 &=\frac{\sigma^2}{N}\,[1+(m-1)\rho].
\end{align}
相对于独立观测，\term{设计效应}{Design Effect}{}为 $1+(m-1)\rho$，等效样本量为 $\frac{N}{[1+(m-1)\rho]}$。例如，每个问题生成五个答案且 $\rho=0.5$，一万个观测的均值方差相当于约 $3333$ 个独立观测。此式依赖等组大小、等相关及组间独立假设；它说明重复扩增增加的存储与训练成本，与其带来的信息增量不成比例。

\subsection{分类维度}
存储格式、领域、记录结构与监督类型是可组合的属性。例如，一个 JSONL 文件可以同时包含数学领域、多轮消息和偏好样本。数据分类可按以下六个维度描述。

\begin{enumerate}
\item \emph{物理格式}描述字节如何编码和组织，例如纯文本、JSONL、CSV、Parquet、Arrow IPC 或 TAR。格式影响压缩、随机访问和扫描成本，不定义训练目标。
\item \emph{逻辑样本模式}描述记录的语义结构，例如文档、平行句对、指令与回答、消息序列、偏好对、强化学习提示或推理轨迹。
\item \emph{消费目标}描述如何使用记录，例如分词器训练、预训练、监督微调、奖励学习、偏好优化、带可验证奖励的强化学习、蒸馏或评估。
\item \emph{内容与能力}描述语言、领域以及代码、数学、工具使用、安全等能力属性。一条记录可以具有多个标签；标签也可能带有置信度。
\item \emph{模态}描述文本、图像、音频和视频等内容，以及跨模态对应关系。时间戳、图像区域和媒体引用属于语义的一部分。
\item \emph{治理与血缘}描述来源快照、使用条件、允许用途、敏感信息、保留期限、质量结论和删除关系。
\end{enumerate}

\term{思维链}{Chain of Thought}{CoT}作为多步推理轨迹的逻辑呈现，属于样本的内容结构属性，其损失函数掩码方式需依据下游消费目标专门设定。例如，同一份推理轨迹数据可派生为仅在最终答案位置计算损失的弱监督视图、对思考过程全量监督的预训练或微调视图，或仅提取初始输入用于强化学习试错生成的提示词视图。当消费目标发生变化时，应构建带独立版本号的物化视图，避免直接原地修改底层基础记录。

\subsection{数据记录格式}
\term{数据模式}{Schema}{}定义字段类型、必要关系和约束。统一数据系统宜采用“公共信封加类型化载荷”：公共信封记录逻辑标识、来源、模式版本、语言、模态和治理信息；载荷由明确的类型标记选择文档、消息、偏好等结构。这样的\term{类型化记录}{Typed Record}{}可以同时检查公共约束和任务约束。

例如，一条偏好记录含公共提示、被偏好的回答和另一回答，二者必须在同一提示语境中可比较；当作两条无关联文本分别筛选会使这一关系失效。工具对话需要保存消息顺序、调用标识与返回结果之间的配对。推理记录应区分过程、最终答案和验证器结论；媒体记录应保留资源引用与对齐信息。把所有内容先展平为文本再尝试恢复结构，会丢失这些关系。

\term{模式适配器}{Schema Adapter}{}负责从指定来源版本转换到统一模式。它必须声明输入版本、输出版本、字段映射和失败条件。未知模式进入隔离队列，仅凭几个字段名推测语义便继续训练会引入难以追查的错误。常见指令格式与消息格式即使名称相同，也可能在角色、工具消息或媒体字段上存在版本差异。\footnote{面向模型接口的消息结构与数据仓库的归档结构职责不同。归档保留训练器不消费的治理字段，接口适配层按调用要求选择输出字段。}

记录标识 $\texttt{record\_id}$ 与分组标识 $\texttt{group\_id}$ 也应分离：前者标识逻辑记录，后者表达不能跨越训练与评估边界的家族。不同来源的局部编号可能相同，因此记录定位至少包含来源数据集标识和局部记录标识。重复来源和派生视图形成的是多对多关系，一条可覆盖的“来源字符串”承担不了这一结构。

\section{数据版本构建}

\subsection{数据层与控制层}
\term{数据血缘}{Data Lineage}{}记录对象由哪些父对象经过哪些变换产生。图\ref{fig:data-pipeline}把字节处理与控制信息分开：数据层负责变换内容，控制层负责确定变换为何被允许、采用哪个版本以及何时可发布。

\begin{figure}[htbp]
\centering
\begin{tikzpicture}[>=stealth,every node/.style={font=\small},box/.style={draw,rounded corners,align=center,minimum width=2.5cm,minimum height=1cm}]
\node[box] (raw) at (0,0) {原始快照\\字节与来源};
\node[box] (canonical) at (3.2,0) {规范记录\\类型与治理};
\node[box] (view) at (6.4,0) {消费视图\\词元与标签};
\node[box] (shard) at (9.6,0) {训练分片\\块与跨度};
\draw[->] (raw)--node[above]{解析}(canonical);
\draw[->] (canonical)--node[above]{物化}(view);
\draw[->] (view)--node[above]{组织}(shard);
\node[draw,align=center,minimum width=11.6cm,minimum height=.85cm] (control) at (4.8,-2) {控制层：策略、版本、隔离、检查点、血缘与发布清单};
\foreach \n in {raw,canonical,view,shard} {\draw[->,dashed] (control.north -| \n.south)--(\n.south);}
\end{tikzpicture}
\caption[数据工程分层体系与变换路径]{数据工程分层体系与变换路径。涵盖原始数据、语义记录与训练制品分层。虚线表示控制约束，实线表示内容变换。}\label{fig:data-pipeline}
\end{figure}

原始层保存可追溯的快照，规范层保存经过解析并保持语义的记录，训练层保存针对特定分词器、目标和布局物化的制品。三者各有分工：训练词元通常不足以恢复解析前的字节，清洗文本也未必能够重建媒体、许可记录和原始边界。各层按其用途制定保留与删除策略。

\subsection{确定性与不同种类的哈希}
\term{确定性变换}{Deterministic Transformation}{}要求在输入及显式配置相同的条件下产生相同结果。配置包括代码、模式适配器、规范化规则、词表、特殊词元、随机种子以及所有影响内容的外部模型版本。输出复现需要同时固定种子、输入顺序、随机数算法和并发归并规则。

必须区分原始字节哈希、提取文本哈希、规范记录哈希、分词器输入哈希和去重规范化哈希。它们回答不同的相等性问题。去重时忽略空白，可能适合普通散文，却会破坏代码缩进或表格含义。规范记录的哈希需要固定序列化规则，例如字段排序、编码、空值与数字表示；序列化规则不固定时，语义相同的对象会得到不同摘要。

密码学哈希核对内容完整性，来源授权与清单身份通过相应凭据验证。如果攻击者同时替换内容与摘要，内部一致性仍然成立。可信注册表或签名等机制属于控制层；来源治理需要这些机制，内容指纹本身覆盖不到这一层。

\section{数据质量控制}

\subsection{过滤后的学习分布}
设输入样本分布为 $P(z)$，样本以概率 $a(z)\in[0,1]$ 被接受，则接受后的分布为
\begin{equation}
Q(z)=\frac{a(z)P(z)}{\mathbb E_{P}[a(Z)]},\qquad 0\le a(z)\le1,\quad \mathbb E_P[a(Z)]>0.
\label{eq:data-filter}
\end{equation}
若平均接受率为零，处理结果是空数据集，上述归一化分布在此无定义；发布流程应报告无可用样本并停止该分支。
这表明质量筛选既减少不合格内容，也重新加权语言、领域和表达方式。若质量评分器偏好某一种文体，保留率描述数据筛选结果，目标能力还需在相应任务上测量。阈值应按记录类型和目标用途解释：代码的有效性、对话的角色一致性、偏好对的可比较性与网页的正文密度是不同指标。

解析、规范化、敏感信息处理和质量筛选应具有分开的阶段及原因码。解析失败意味着无法可靠解释结构；质量拒绝意味着结构可解释但不满足某项条件。将所有失败统称为“坏数据”，会使源格式升级与真实质量变化无法区分。统计时还应声明分母：源记录保留率、规范记录保留率和有效词元保留率可能给出截然不同的结论。表\ref{tab:data-pipeline-stages}总结了典型网络预训练语料清洗流水线各阶段的判定规则与典型保留率。

\begin{table}[htbp]
\centering
\caption{典型预训练语料处理规则与保留率}
\label{tab:data-pipeline-stages}
\begin{tabularx}{\linewidth}{lp{28mm}p{38mm}X}
\toprule
处理阶段 & 核心判定规则/模型 & 典型目标与特征 & 典型累积保留率 \\
\midrule
原始抓取 & Warc/HTML 响应提取 & 状态码 200、非空文档 & $100\%$（基准） \\
文本抽取与清洗 & 正文提取（如 Trafilatura） & 去除 HTML 标签、脚本及导航栏 & 约 $70\%\sim 85\%$ \\
浅层启发式过滤 & 字符长度、标点比例、停用词率 & 滤除乱码、SEO 堆叠、词表外字符 & 约 $40\%\sim 60\%$ \\
语言识别 & FastText / CLD3 语言分类器 & 识别主语言并过滤低置信度文档 & 约 $30\%\sim 45\%$ \\
质量分类过滤 & 分类器（如 基于 Wikipedia/教材训练） & 筛选高信息密度与优质行文表达 & 约 $15\%\sim 25\%$ \\
文本去重 & MinHash+LSH / 后缀数组 & 去除网页爬取镜像、样板重复片段 & 约 $8\%\sim 15\%$ \\
安全合规清洗 & PII 屏蔽、毒性与违规内容检测 & 剔除敏感身份信息与不良内容 & 约 $7\%\sim 14\%$ \\
\bottomrule
\end{tabularx}
\end{table}

\term{个人可识别信息}{Personally Identifiable Information}{PII}检测可以结合规则、实体识别和上下文判定，但模式匹配只能覆盖已定义模式。匿名化需要评估多个字段组合后的重识别风险，电话号码等直接标识和可关联的间接标识均需处理。审计日志应保留稳定原因和定位信息，避免把被移除的秘密再次写入普通日志。对敏感记录采用隔离或删除后，还需要同步处理其派生视图。

\subsection{污染与重复}
\term{评估污染}{Evaluation Contamination}{}指评估信息通过训练、调参或数据选择进入被评估系统，其结果因而失去了独立泛化的解释力。问题全文、答案、解释、翻译和近似改写都可能构成信息通道。仅检查整条记录的哈希相等，只排除了很窄的一类重叠。

设评估题的片段集合为 $A$，训练文档为 $B$。当短题被嵌入长文档时，Jaccard 相似度可能很低，而\term{包含率}{Containment}{} $C(A,B)=\frac{|A\cap B|}{|A|}$ 仍为 $1$。因此污染检查需要与任务匹配的粒度：整记录、片段、问题家族以及答案相关内容。包含率与相似度的区分见 Broder 的原始研究\cite{broder1997resemblance}。

污染检查的覆盖范围由检测方法、阈值和参考集合决定，报告应说明这些条件。保留评估参考集的访问边界也很重要：为了检测而把答案完整写入通用训练日志或生成提示，会把污染扩散到新的位置。数据去重研究显示训练重复与记忆、训练评估重叠存在值得控制的联系\cite{lee2022deduplicating}。

\section{数据去重}

\subsection{精确相等的判定对象}
\term{精确去重}{Exact Deduplication}{}先对指定对象进行确定性规范化，再比较其摘要；需要排除碰撞时进一步比较规范化字节。对象可以是文档、对话或完整偏好对。对结构化记录只比较回答文本，会把语境不同的合法样本误判为相同。

重复记录的删除应保留所有来源别名与治理关系。一个内容在多个位置出现，这些来源的使用条件仍各自独立。代表记录的选择需明确优先级，例如允许用途、质量、来源可信度、时间与稳定标识的顺序。并发完成顺序一旦决定保留者，同一输入就会产生不同训练版本。常见去重算法的机制、复杂度及适用场景如表\ref{tab:dedup-comparison}所示。

\begin{table}[htbp]
\centering
\caption{主流文本去重算法机制与适用性对比}
\label{tab:dedup-comparison}
\begin{tabularx}{\linewidth}{lp{28mm}p{25mm}X}
\toprule
去重方法 & 匹配机制 & 计算/空间复杂度 & 特点与工程适用场景 \\
\midrule
精确哈希 (MD5/SHA256) & 规范化字符串字节全等 & $\mathcal{O}(N)$ 时间 / $\mathcal{O}(N)$ 空间 & 识别完全相同文档，无法处理包含轻微修改或不同页眉页脚的文档 \\
MinHash + LSH & $q$-gram Jaccard 相似度 & $\mathcal{O}(N \cdot m)$ 时间 / 倒排索引空间 & 高效检出高相似度文档（如 Jaccard $\ge 0.8$），适合海量 Web 网页清洗 \\
后缀数组 (Suffix Array) & 连续子串匹配（最长公共前缀） & $\mathcal{O}(N \log N)$ 构造 / $\mathcal{O}(N)$ 空间 & 跨文档级细粒度长重复词元检测与消除，有效缓解模型退化与记忆 \\
语义向量相似度 & 嵌入向量最近邻（ANN/HNSW） & $\mathcal{O}(N \cdot d)$ 编码 / 索引高开销 & 识别改写、多语言等意译重复，但算力开销极大，多用于精细后训练数据 \\
\bottomrule
\end{tabularx}
\end{table}

\subsection{MinHash 去重}
把规范文本切成长度为 $q$ 的连续片段，称为\term{连续片段}{Shingle}{}，得到集合 $A$。这里采用集合而非多重集，重复片段出现次数不增加权重。两条记录的\term{Jaccard 相似度}{Jaccard Similarity}{}为
\begin{equation}
J(A,B)=\frac{|A\cap B|}{|A\cup B|}.
\end{equation}
$q$ 太小会使常用表达大量重合，太大则对局部改写敏感。空集合需要另行约定，例如送入短文本专门处理；$\frac{0}{0}$ 在该式中没有定义。

\begin{assumption}[最小哈希的理想化模型]
片段来自有限全集，$A\ne\varnothing$ 且 $B\ne\varnothing$；空集合已由前述短文本流程单独处理。每次独立地对全集作均匀随机排列 $\pi$，并以排列次序最小的片段作为集合签名。推导暂不考虑有限哈希值的碰撞及非理想哈希族偏差。
\end{assumption}

\term{最小哈希}{MinHash}{}用 $h_\pi(A)=\argmin_{x\in A}\pi(x)$ 表示集合。两个最小值相同，当且仅当 $A\cup B$ 中排列最靠前的元素落在 $A\cap B$ 中。并集内每个元素成为最小者的概率相等，因此
\begin{equation}
\Pr[h_\pi(A)=h_\pi(B)]=\frac{|A\cap B|}{|A\cup B|}=J(A,B).
\label{eq:data-minhash}
\end{equation}
该关系把集合相似度转化为签名碰撞概率\cite{broder1997resemblance}。采用 $m$ 次独立排列，令 $I_j$ 为第 $j$ 个签名相等的指示量，则
\begin{equation}
 \begin{aligned}
  \widehat J &= \frac{1}{m}\sum_{j=1}^{m}I_j, \\
  \mathbb E[\widehat J] &= J, \\
  \operatorname{Var}(\widehat J) &= \frac{J(1-J)}{m}.
 \end{aligned}
\end{equation}
例如 $J=0.8,m=100$ 时标准差为 $0.04$。该标准差描述随机估计的离散程度；单次估计的误差范围需结合分布假设和置信水平计算。工程上用哈希函数近似随机排列，函数相关性、摘要截断和碰撞都会改变上述理想性质。

\subsection{局部敏感哈希候选集}
对 $N$ 条记录逐对计算相似度需要 $O(N^2)$ 次比较。\term{局部敏感哈希}{Locality-Sensitive Hashing}{LSH}把签名分成 $b$ 个带，每带 $r$ 行，满足 $m=br$；同一带内所有行相同的记录进入同一候选桶。若两记录相似度为 $s$，在独立行假设下，一个带完全相同的概率为 $s^r$，所有带均不同的概率为 $(1-s^r)^b$，所以成为候选的概率为
\begin{equation}
P_{\mathrm{cand}}(s)=1-(1-s^r)^b.
\label{eq:data-lsh}
\end{equation}
该函数在 $s\approx (1/b)^{1/r}$ 处产生 S 型曲线，可控制召回与候选规模。

\begin{example}[带划分参数对候选概率的影响]
设签名长度 $m=100$。方案一取 $b=20,r=5$；方案二取 $b=10,r=10$。当相似度 $s=0.8$ 时，
方案一候选概率为 $1-(1-0.8^5)^{20}\approx0.999$，
方案二为 $1-(1-0.8^{10})^{10}\approx0.669$。
较小的 $r$ 增大召回率，但也引入更多低相似度候选对，增加后验精确验证的开销。
\end{example}

\begin{figure}[htbp]
\centering
\begin{tikzpicture}[>=stealth,every node/.style={font=\small},b/.style={draw,align=center,minimum height=.9cm}]
\node[b] (a) at (0,0) {片段集合};
\node[b] (b) at (2.8,0) {$m=br$ 个签名};
\node[b] (c) at (6,0) {带内全等\\带间取并集};
\node[b] (d) at (9.3,0) {候选对\\真实相似度确认};
\draw[->] (a)--(b);\draw[->](b)--(c);\draw[->](c)--(d);
\node[align=center] at (5,-1.35) {候选召回控制计算量；确认阈值控制候选中的误判};
\end{tikzpicture}
\caption[近似去重候选生成与精确确认流程]{近似去重候选生成与精确确认流程。候选生成与精确确认承担不同职责。}\label{fig:data-lsh}
\end{figure}

对候选计算真实 Jaccard 可以排除候选中的误报，从未进入候选的漏报则无法由此找回。候选桶中有 $u$ 条记录时，完全展开仍产生 $\frac{u(u-1)}{2}$ 对。常见模板会形成热点桶，限制桶大小虽能保护资源，却改变召回率；这种限制属于策略版本的一部分，隐藏成实现细节会掩盖它对结果的影响。

\subsection{近重关系的非传递性}
令 $A=\{1,2\}$，$B=\{1,2,3\}$，$C=\{2,3\}$，阈值取 $\frac{2}{3}$。此时 $J(A,B)=J(B,C)=\frac{2}{3}$，但 $J(A,C)=\frac{1}{3}$。如果把近重边的连通分量全部只留一条，就可能删除与代表记录并不相似的内容。

分组和删除需要不同语义。为了防止泄漏，可以保守地把整个关联分量放入同一个划分；为了删除内容，则可按稳定优先级扫描，只有与已经保留的某个代表直接达到阈值的记录才删除。该策略不以求最小代表集合为目标，也可能保留多个彼此通过中间节点相关的记录；它换取的是可解释的直接删除依据。

\begin{algorithm}[稳定代表去重]
\label{alg:data-dedup}
\AlgoInput{类型一致的记录、固定规范化规则、相似度阈值 $\tau$、候选索引及稳定全序。}
\AlgoOutput{保留记录与删除依据。}
\AlgoState{保留集合 $R$，重复映射 $D$。}
\begin{enumerate}
\item 按治理资格、质量和稳定标识确定全序；初始化 $R=D=\varnothing$。
\item 依次读取记录 $x$，从候选索引取得已在 $R$ 中的候选，按同一全序排列。
\item 对每个候选 $y$ 计算实际相似度；若有 $J(x,y)\ge\tau$，记录直接依据 $(x,y,J)$，并保留 $x$ 的所有来源血缘。
\item 若没有确认匹配，将 $x$ 加入 $R$。处理完所有记录后停止。
\end{enumerate}
\textbf{不变量：}每条因近重被删除的记录都有一个直接达标且仍被保留的代表。候选漏检仍可能导致重复保留，必须单独评估召回。
\end{algorithm}

\section{训练分布的构造}

\subsection{分组哈希与时间边界}
\term{分组划分}{Group Split}{}要求同一个泄漏家族只进入训练、验证或测试中的一个集合。设稳定组标识为 $g$，固定盐值为 $s_0$，把哈希映射到 $u(g)\in[0,1)$。例如 $u<0.8$ 分配到训练，$0.8\le u<0.9$ 分配到验证，其余进入测试。同一组的新增记录仍进入同一集合；与按当前行号洗牌相比，它更适合增量数据。

这里稳定的是组到集合的映射，不是精确比例。组大小不等时，样本比例和词元比例都可能偏离 $80:10:10$。哈希还必须与目标变量及来源顺序没有系统性关联。盐值在划分前固定，以保持数据分配可复现。

如果新增近重边把两个已经落入不同集合的组连接起来，原划分的隔离条件已经失效。保持所有旧分配与严格分组隔离在此无法兼得；需要隔离新增冲突、重建划分或发布新的评估版本。面向未来数据的应用还应采用时间边界，确保训练时不可见的信息没有进入特征构造。分词器学习、统计归一化和基于语料的阈值拟合也属于需要限定训练范围的过程。

\subsection{温度混合与目标校正}
设清洗后有 $D$ 个合格来源，第 $i$ 个来源具有 $n_i>0$ 个采样单位。\term{温度混合}{Temperature-Based Mixing}{}常写为
\begin{equation}
p_i=\frac{n_i^\alpha}{\sum_{j=1}^{D}n_j^\alpha},\qquad 0\le\alpha\le1.
\label{eq:data-mixing}
\end{equation}
$\alpha=1$ 保持自然数量比例，$\alpha=0$ 等概率选择来源，较小的 $\alpha$ 提高小来源的相对权重。式中的 $n_i$ 必须声明是文档数、词元数还是其他单位，同一配置里混用不同单位会使配比失去意义。

例如两个来源分别有 $900$ 和 $100$ 条记录，取 $\alpha=\frac{1}{2}$，则 $p=(0.75,0.25)$。若有放回抽取 $1200$ 次，期望分别抽取 $900$ 和 $300$ 条，单条记录平均使用次数分别为 $1$ 和 $3$。若只对各来源各读一遍并拼接，实际记录比例仍为 $0.9:0.1$，并未实现指定目标。

假设两个来源平均长度分别为 $100$ 和 $400$ 个有效词元，按上述文档概率采样，则长期词元份额近似为
\begin{equation}
\widetilde p_i=\frac{p_i\mathbb E[L_i]}{\sum_jp_j\mathbb E[L_j]},
\end{equation}
得到 $\widetilde p=(\frac{3}{7},\frac{4}{7})$。因此目标配比、抽中文档配比、物化词元配比和最终有效标签配比必须分别统计。截断、尾部丢弃和标签屏蔽会使它们进一步不同。

若目标分布的来源权重为 $q_i$，实际按 $p_i$ 抽样，可以采用\term{重要性采样}{Importance Sampling}{}校正：对来自 $i$ 的损失乘以 $\frac{q_i}{p_i}$。因为
\begin{equation}
\sum_i p_i\,\mathbb E_i\!\left[\frac{q_i}{p_i}\ell_\theta\right]
=\sum_iq_i\mathbb E_i[\ell_\theta],
\end{equation}
期望目标恢复为指定加权风险。前提是 $q_i>0$ 时 $p_i>0$，并且期望存在；极小的 $p_i$ 会产生大权重和高方差。若本来就希望改变学习目标，再施加该校正会抵消这一意图。

\term{课程学习}{Curriculum Learning}{}让 $p_i$ 随训练步变化。它改变参数经历的样本顺序与梯度路径，最终累计比例描述不到这一过程。\term{合成数据}{Synthetic Data}{}同样需要来源问题家族、生成模型、提示、解码配置和验证器版本。生成器自评筛选可能放大其偏好，验证器通过也只证明所检验的条件；评估集应与生成、筛选和阈值调整隔离。

\section{训练样本的物化}

\subsection{训练样本消费视图}
\term{词元物化}{Token Materialization}{}把规范记录转换为固定词元与监督信号。它需要冻结分词器、特殊词元、对话模板、截断规则、标签策略和位置处理。仅记录词表名称无法唯一确定结果。对于偏好对，共同提示的处理必须一致；对于推理样本，应显式决定过程是否参与损失。

训练视图至少定义输入词元、有效预测位置和必要的边界信息。自回归目标预测的是下一词元，输入与目标的错位必须与样本边界一致。有效标签数由监督掩码决定：提示词元、填充、边界和被屏蔽的内容都可能不计入损失。从文档到批次张量的数据物化与组批流水线如图\ref{fig:pretraining-data-modules}所示。

\begin{figure}[htbp]
\centering
\includegraphics[width=0.92\textwidth]{\subfix{../../figures/theory/pretraining-data-modules.pdf}}
\caption{训练批次数据物化与组批流水线}
\label{fig:pretraining-data-modules}
\end{figure}

\subsection{连续流与不可拆样本的容量公式}
\term{序列打包}{Sequence Packing}{}把变长内容组织为长度上限为 $T$ 的模型输入块。首先区分允许跨块切分的连续流与不允许切分的原子样本。令 $L_i$ 已包含规则要求的分隔符或边界词元，总长度 $M=\sum_iL_i$。

对连续流，若尾部不足一块就丢弃，则
\begin{equation}
 \begin{aligned}
  K_{\mathrm{drop}} &= \left\lfloor \frac{M}{T}\right\rfloor, \\
  R &= M-K_{\mathrm{drop}}T, \qquad 0\le R<T.
 \end{aligned}
\end{equation}
若保留尾部并填充，则对于 $M>0$，
\begin{equation}
 \begin{aligned}
  K_{\mathrm{pad}} &= \left\lceil \frac{M}{T}\right\rceil, \\
  P &= K_{\mathrm{pad}}T-M.
 \end{aligned}
\end{equation}
空输入约定零块。块数由输入长度和块容量确定，填充量由总容量减去输入长度得到。

如果每条样本不可拆分，$L_i\le T$，则 $\lceil \frac{M}{T}\rceil$ 只是块数下界，实际值由装箱策略决定。例如三个长度均为 $6$ 的样本，$T=10$，总长 $18$ 给出下界 $2$，但任意两条都无法同块，必须使用 $3$ 块。超过 $T$ 的样本还需要明确选择截断、切分或隔离，留给隐式溢出行为会使边界不可见。

长度为 $(5,4,3)$、$T=8$ 时，按顺序采用放不下便开启下一块的策略，得到第一块长度 $5$、第二块长度 $7$，共填充 $4$ 个位置，内容利用率为 $\frac{12}{16}=75\%$。连续流丢尾则生成一块、丢弃四个词元，已生成块利用率为 $100\%$，但数据损失率为 $\frac{1}{3}$。只汇报利用率会掩盖这一差别。特殊词元占用率与有效标签占用率还应另外统计。

\begin{figure}[htbp]
\centering
\begin{tikzpicture}[x=.72cm,y=.65cm,every node/.style={font=\small}]
\foreach \i in {0,...,7} {\draw (\i,0) rectangle +(1,1);\draw (\i,-2) rectangle +(1,1);}
\node at (2.5,.5) {样本 A：5};\node at (6.5,.5) {填充：3};
\node at (2,-1.5) {样本 B：4};\node at (5.5,-1.5) {样本 C：3};\node at (7.5,-1.5) {空};
\draw[very thick] (5,0)--(5,1);\draw[very thick] (4,-2)--(4,-1);\draw[very thick] (7,-2)--(7,-1);
\node[anchor=west] at (8.4,.5) {块 1：$[0,5)$ 属于 A};
\node[anchor=west] at (8.4,-1.5) {块 2：$[0,4)$ 属于 B};
\end{tikzpicture}
\caption[不可拆样本顺序打包示意]{不可拆样本顺序打包示意。每格为一个位置，粗线是来源边界；是否允许跨边界注意力仍由掩码策略决定。}\label{fig:data-packing}
\end{figure}

\subsection{样本隔离约束}
文档装块与注意力可见性分别配置。连续预训练流可以显式允许跨文档条件关系；若目标要求样本独立，则需要块内分段因果掩码。令 $g(t)$ 为位置 $t$ 所属样本，允许位置 $t$ 读取 $j$ 的条件为 $j\le t$ 且 $g(j)=g(t)$，填充位置另行处理。\bookxref{ch:attention}讨论了掩码归一化，包括全屏蔽行的边界。

仅插入结束词元并不阻断注意力。标签也必须隔离：若前一样本的末位置被要求预测下一样本的首词元，即使注意力已隔离，损失仍引入跨样本预测。可以在每个样本加入开始标记，由其预测本样本首词元，或屏蔽没有合法本样本前缀的预测位置。位置编号是否重置应写入布局规范。\footnote{打包后的注意力掩码在概念上可以是分块矩阵；物理实现可以使用段边界等紧凑表示。段边界可用于构造或执行相应的分块掩码。}

\keyconcept{序列打包（Packing）若仅拼接结束词元而未构建分段因果注意力掩码，模型将在自注意力计算中跨越独立文档边界读取无关上下文；同时末尾位置若未屏蔽下一样本首词元的预测损失，将引入虚假的跨样本因果关联。}

\begin{warning}[序列打包的隐式上下文污染]
在微调或需要严谨推理的语料上进行序列打包时，若底层算子（如部分 FlashAttention 实现）未正确配置文档切分累积长度（\texttt{cu\_seqlens}），注意力矩阵将退化为普通因果掩码，导致后一样本直接“窥见”前一样本内容，引发训练损失虚低与推理幻觉。
\end{warning}

\begin{algorithm}[保持来源的原子样本打包]
\AlgoInput{按确定顺序排列的已序列化样本，最大长度 $T$，超长与尾部策略。}
\AlgoOutput{词元块、标签掩码、样本段及来源跨度。}
\AlgoState{当前块 $B$、已用长度 $u$、跨度表 $S$。}
\begin{enumerate}
\item 初始化空块。读取下一样本，先执行已声明的超长策略，使进入打包器的原子样本长度不超过 $T$。
\item 若 $u+L_i>T$，按尾部策略封存当前块，记录填充或丢弃量；重新建立空块。
\item 追加样本，记录半开区间 $[u,u+L_i)$、来源记录与源词元偏移，并设置段标识和合法标签；更新 $u$。
\item 输入耗尽时按同一尾部策略处理最后一块，停止。
\end{enumerate}
\textbf{不变量：}块长不超过 $T$；每个非填充位置能定位到来源或明确标记为插入词元；样本边界与标签策略一致。按顺序一次遍历的组织成本为 $O(\sum_iL_i)$，不含分词成本。
\end{algorithm}

\subsection{来源跨度与分片边界}
\term{来源跨度}{Provenance Span}{}将块中的半开区间映射到来源数据集、记录、源偏移以及变换版本。填充应有空来源，插入的特殊词元应注明生成规则，伪装为原文词元会污染来源统计。跨度可支持训练贡献统计、错误定位与删除影响分析；只保存“本块来自三个文档”不足以定位具体内容。

\term{分片}{Shard}{}是存储或调度单位，\term{训练块}{Training Block}{}是模型输入单位，两者属于不同的抽象层次。一个分片可以容纳多个完整块，也可以在另一种明确格式中保存尚未打包的文档。若规定分片只含完整 $T$ 长块，容量上限为 $S$ 个位置，则最多容纳 $\lfloor \frac{S}{T}\rfloor$ 块，且要求 $S\ge T$。例如 $S=18,T=8$，只能放两块即 $16$ 个位置；余下两个位置不能容纳第三块。文件结束也不是文档结束或训练序列结束。

下面将来源采样、序列打包与标签统计合在一起计算，比较不同计量口径。

\begin{example}[从来源采样到有效监督词元]
考虑两个通过治理筛选的来源，共 $900$ 条短记录与 $100$ 条长记录。采用式\eqref{eq:data-mixing}中的 $\alpha=\frac{1}{2}$，按文档有放回抽样 $1200$ 次，期望使用次数为 $900$ 与 $300$。若序列化后长度恒为 $100$ 与 $400$，则期望物化长度为 $210000$ 个词元；其中长来源贡献 $120000$，即 $\frac{4}{7}$。

现在取一个具体已经实现该次数的采样序列，按连续流以 $T=1024$ 打包。保留尾部时有 $\lceil\frac{210000}{1024}\rceil=206$ 块，容量为 $210944$，填充 $944$ 个位置。丢尾时有 $205$ 块，保留 $209920$ 个词元，丢弃 $80$ 个。由于尾部来自具体排序中的记录，其来源份额不一定仍为 $\frac{4}{7}$。若改成不可拆记录，这两个块数公式便不再直接适用。

进一步假设标签策略屏蔽所有提示，共屏蔽 $42000$ 个实际词元；在保尾条件下有效标签最多为 $168000$，还需扣除策略要求屏蔽的边界位置。按块容量统计的有效标签率最多为 $\frac{168000}{210944}\approx79.64\%$，低于内容利用率约 $99.55\%$。

来源记录数、抽样次数、实际词元数和有效标签数分别描述数据处理的不同阶段，应在各阶段独立统计。
\end{example}

\section{数据流水线的运行管理}

\subsection{读取状态与数据消费}
\term{流式读取}{Streaming}{}降低预先装载需求，读取速度仍受存储与传输限制。内存映射依赖页缓存、存储吞吐和访问局部性。有限缓冲区洗牌只在局部窗口内混合，可能保留原始顺序中的长程相关性。若调度器总选择先完成的数据，较慢来源的实际占比还可能低于目标混合比例。

恢复点必须覆盖分片位置、记录偏移、随机数状态、洗牌缓冲、未封存词元以及训练器已经提交的进度。只恢复文件游标，可能重复或跳过缓冲区中的样本。改变工作进程数或分片分配方式，也可能改变顺序与尾部结果。\term{恰好一次}{Exactly-Once}{}在这里指逻辑训练事件的提交语义，与文件是否只被读取一次无关；是否重复更新参数，取决于数据消费状态与训练检查点如何一致提交。

遇到坏记录静默跳过会改变分布，重试直到成功又可能阻塞整个输入。正确做法是固定失败分类和隔离策略，计量其来源分布，并把最终纳入或排除的决定记录到版本中。批处理适合生成不可变制品，增量处理适合吸收新增快照；二者都需要相同的语义规则。

\subsection{幂等发布与版本身份}
\term{幂等性}{Idempotency}{}表示同一逻辑构建重复提交不会产生不同的可见结果。构建标识可以由源快照集合、模式与适配器、代码和策略版本、分词与布局配置的规范表示计算得到。它表示构建规格；最终内容摘要表示构建产物，两者应分开，以避免让清单把自己的摘要纳入计算而形成自引用。

\term{发布清单}{Manifest}{}列出制品、摘要、记录与词元计数、来源范围、策略版本和必要血缘。中间文件先写入隔离的暂存空间，完整制品准备后才发布清单或更新可见指针。单一文件系统中的原子重命名与对象存储中的条件指针更新具有不同条件；把多个对象依次上传称为整体原子发布，与该词的含义不符。

\begin{algorithm}[可重试的数据版本发布]
\AlgoInput{不可变源快照与完整构建规格。}
\AlgoOutput{可消费版本标识。}
\AlgoState{构建标识、暂存制品列表、阶段完成记录。}
\begin{enumerate}
\item 计算构建标识；若已有已发布版本，比较规格与产物登记，匹配则复用，否则报告一致性冲突。
\item 按阶段生成不可变制品，记录父对象与完成状态。重试只复用经过完整性确认的阶段结果，完成状态由内容完整性与阶段提交记录共同确认。
\item 生成独立内容摘要和清单，记录成功、拒绝、隔离、派生及有效标签数量；所有必需对象就绪后进入发布步骤。
\item 以平台支持的原子提交或条件更新暴露最终指针。并发提交者若发现版本已经发布，核对后复用；同一规格产生不同内容时停止发布并保留诊断依据。
\end{enumerate}
\textbf{不变量：}消费者只能发现完整版本；同一版本标识不覆盖为另一内容。失败发生在发布前时可重试，发布后则返回已存在的版本。
\end{algorithm}

内容摘要可以发现字节变化，但如果攻击者能同时替换制品与清单，摘要核验还需绑定可信发布身份。发布目录应把认可的摘要保存在独立信任边界内，或者由受控签名及验证密钥建立真实性。公开清单通常记录聚合数与受控审计资产的引用，不直接暴露隔离原文、个人标识或推理轨迹；细粒度血缘的可追溯性必须与访问权限同时设计。

对于一进一出的阶段，输入记录数应等于成功、拒绝与隔离数之和；对于一条文档切成多块或多条记录合并的阶段，简单行数守恒不再成立，必须使用父子映射和词元账目。并发执行时，完成集合需逐项记录，最大处理编号仅描述进度上界。

\subsection{删除传播与模型影响}
\term{删除血缘}{Deletion Lineage}{}用于从撤回的源记录定位规范记录、消费视图、词元跨度、分片以及使用过这些制品的训练任务。精确去重后如果删除了来源别名，这条链就会断裂；代表内容与全部来源别名共同纳入治理。

删除请求首先在流式存储与分片调度队列中拦截未来消费，并根据数据治理规范触发受影响数据版本的重新构建或软删除撤销。对于已完成的模型检查点，数据撤回并不等价于参数权重的自动消除：存储层删除、下游数据制品重构与模型参数遗忘（Machine Unlearning）对应不同的工程与算法处置机制。\footnote{用于防止再次导入的删除标记也可能含有可关联标识，其保留形式、权限和期限应单独设计。}

\begin{note}[数据版本的追溯记录]
每个训练数据版本应保存获准使用的来源、内容与标签的变换记录，以及训练器的实际消费进度。来源失效时，可沿这些记录定位受影响的数据快照、训练任务与模型产物，并安排后续处理。
\end{note}

数据审计还应分别记录去重候选召回、隔离率、来源跨度覆盖率、重试产生的逻辑重复量和版本完整性，用于定位各处理阶段的问题。

\chapterreferences
\myendchapterdocument
